\documentclass[a4paper,12pt]{article}

\usepackage{fontspec}
\usepackage{polyglossia}
\usepackage{amsmath}
\usepackage{amssymb}
\usepackage{geometry}

\setmainlanguage{russian}
\setmainfont{CMU Serif}
\newfontfamily\cyrillicfont{CMU Serif}

\geometry{
    left=25mm,
    right=15mm,
    top=20mm,
    bottom=20mm
}

\begin{document}

\begin{center}
    \Large\textbf{Лабораторная работа 11}\\
    \large\textbf{Прямое и обратное распространение}\\
    \vspace{0.5cm}
    28 апреля 2026 г.
\end{center}

\section*{Цель работы}

Цель работы --- найти выражения для производных ошибки по входным сигналам $x_1$ и $x_2$ в простой нейронной сети с использованием правила обратного распространения ошибки.

\section*{Описание нейронной сети}

На вход сети подаются два скалярных значения:

\[
x_1, \quad x_2
\]

Каждый вход проходит через сигмоидную функцию:

\[
i_1 = \sigma(x_1)
\]

\[
i_2 = \sigma(x_2)
\]

Сигмоидная функция имеет вид:

\[
\sigma(x) = \frac{1}{1 + e^{-x}}
\]

Производная сигмоидной функции:

\[
\sigma'(x) = \sigma(x)(1 - \sigma(x))
\]

Далее выходы сигмоидного слоя проходят через слои $min$ и $max$:

\[
o_{min} = \min(i_1, i_2)
\]

\[
o_{max} = \max(i_1, i_2)
\]

На последнем слое значения складываются:

\[
y = o_{min} + o_{max}
\]

Необходимо выразить:

\[
\frac{\partial E}{\partial x_1}
\quad \text{и} \quad
\frac{\partial E}{\partial x_2}
\]

через ошибку, которая приходит с выхода сети:

\[
\frac{\partial E}{\partial y}
\]

\section*{Упрощение выражения для выхода}

Для двух чисел сумма минимального и максимального значения равна сумме этих чисел:

\[
\min(i_1, i_2) + \max(i_1, i_2) = i_1 + i_2
\]

Следовательно, выход сети можно записать так:

\[
y = i_1 + i_2
\]

Подставим значения $i_1$ и $i_2$:

\[
y = \sigma(x_1) + \sigma(x_2)
\]

\section*{Производная по первому входу}

Пусть ошибка, которая приходит с выхода сети, равна:

\[
\frac{\partial E}{\partial y}
\]

Найдем производную ошибки по входу $x_1$.

Используем правило цепочки:

\[
\frac{\partial E}{\partial x_1}
=
\frac{\partial E}{\partial y}
\cdot
\frac{\partial y}{\partial x_1}
\]

Так как:

\[
y = \sigma(x_1) + \sigma(x_2)
\]

то производная по $x_1$ равна:

\[
\frac{\partial y}{\partial x_1}
=
\sigma'(x_1)
\]

Подставим производную сигмоиды:

\[
\frac{\partial y}{\partial x_1}
=
\sigma(x_1)(1 - \sigma(x_1))
\]

Тогда:

\[
\frac{\partial E}{\partial x_1}
=
\frac{\partial E}{\partial y}
\cdot
\sigma(x_1)(1 - \sigma(x_1))
\]

\section*{Производная по второму входу}

Аналогично найдем производную ошибки по входу $x_2$.

Используем правило цепочки:

\[
\frac{\partial E}{\partial x_2}
=
\frac{\partial E}{\partial y}
\cdot
\frac{\partial y}{\partial x_2}
\]

Так как:

\[
y = \sigma(x_1) + \sigma(x_2)
\]

то производная по $x_2$ равна:

\[
\frac{\partial y}{\partial x_2}
=
\sigma'(x_2)
\]

Подставим производную сигмоиды:

\[
\frac{\partial y}{\partial x_2}
=
\sigma(x_2)(1 - \sigma(x_2))
\]

Тогда:

\[
\frac{\partial E}{\partial x_2}
=
\frac{\partial E}{\partial y}
\cdot
\sigma(x_2)(1 - \sigma(x_2))
\]

\section*{Итоговые формулы}

Итоговая формула для первого входа:

\[
\boxed{
\frac{\partial E}{\partial x_1}
=
\frac{\partial E}{\partial y}
\cdot
\sigma(x_1)(1 - \sigma(x_1))
}
\]

Итоговая формула для второго входа:

\[
\boxed{
\frac{\partial E}{\partial x_2}
=
\frac{\partial E}{\partial y}
\cdot
\sigma(x_2)(1 - \sigma(x_2))
}
\]

\section*{Вывод}

В ходе работы было выполнено обратное распространение ошибки через простую нейронную сеть.

Входные значения $x_1$ и $x_2$ сначала проходят через сигмоидную функцию, после чего попадают в слои $min$ и $max$.

Так как сумма значений $min$ и $max$ для двух чисел равна сумме этих чисел, выход сети можно упростить:

\[
y = \sigma(x_1) + \sigma(x_2)
\]

Поэтому производные ошибки по входам $x_1$ и $x_2$ выражаются через производную сигмоидной функции и ошибку, приходящую с выхода сети.

\end{document}